\documentclass[10pt,a4paper]{amsart}
\usepackage{amsmath,amssymb,amsthm,graphicx}
\usepackage{verbatim}
\newenvironment{reference}{\comment}{\endcomment}
%\newenvironment{reference}{}{}
\newenvironment{slogan}{\comment}{\endcomment}
\newenvironment{history}{\comment}{\endcomment}
% For commutative diagrams we use Xy-pic
\usepackage[all]{xy}

% We use 2cell for 2-commutative diagrams.
\xyoption{2cell}
% Theorem environments.
%
\theoremstyle{plain}
\newtheorem{theorem}[subsection]{Theorem}
\newtheorem{proposition}[subsection]{Proposition}
\newtheorem{lemma}[subsection]{Lemma}

\theoremstyle{definition}
\newtheorem{definition}[subsection]{Definition}
\newtheorem{example}[subsection]{Example}
\newtheorem{exercise}[subsection]{Exercise}
\newtheorem{situation}[subsection]{Situation}

\theoremstyle{remark}
\newtheorem{remark}[subsection]{Remark}
\newtheorem{remarks}[subsection]{Remarks}

\numberwithin{equation}{subsection}

% Macros
%
\def\lim{\mathop{\mathrm{lim}}\nolimits}
\def\colim{\mathop{\mathrm{colim}}\nolimits}
\def\Spec{\mathop{\mathrm{Spec}}}
\def\Hom{\mathop{\mathrm{Hom}}\nolimits}
\def\Ext{\mathop{\mathrm{Ext}}\nolimits}
\def\SheafHom{\mathop{\mathcal{H}\!\mathit{om}}\nolimits}
\def\SheafExt{\mathop{\mathcal{E}\!\mathit{xt}}\nolimits}
\def\Sch{\mathit{Sch}}
\def\Mor{\mathop{\mathrm{Mor}}\nolimits}
\def\Ob{\mathop{\mathrm{Ob}}\nolimits}
\def\Sh{\mathop{\mathit{Sh}}\nolimits}
\def\NL{\mathop{N\!L}\nolimits}
\def\CH{\mathop{\mathrm{CH}}\nolimits}
\def\proetale{{pro\text{-}\acute{e}tale}}
\def\etale{{\acute{e}tale}}
\def\QCoh{\mathit{QCoh}}
\def\Ker{\mathop{\mathrm{Ker}}}
\def\Im{\mathop{\mathrm{Im}}}
\def\Coker{\mathop{\mathrm{Coker}}}
\def\Coim{\mathop{\mathrm{Coim}}}

% Boxtimes
%
\DeclareMathSymbol{\boxtimes}{\mathbin}{AMSa}{"02}

%
% Macros for moduli stacks/spaces
%
\def\QCohstack{\mathcal{QC}\!\mathit{oh}}
\def\Cohstack{\mathcal{C}\!\mathit{oh}}
\def\Spacesstack{\mathcal{S}\!\mathit{paces}}
\def\Quotfunctor{\mathrm{Quot}}
\def\Hilbfunctor{\mathrm{Hilb}}
\def\Curvesstack{\mathcal{C}\!\mathit{urves}}
\def\Polarizedstack{\mathcal{P}\!\mathit{olarized}}
\def\Complexesstack{\mathcal{C}\!\mathit{omplexes}}
% \Pic is the operator that assigns to X its picard group, usage \Pic(X)
% \Picardstack_{X/B} denotes the Picard stack of X over B
% \Picardfunctor_{X/B} denotes the Picard functor of X over B
\def\Pic{\mathop{\mathrm{Pic}}\nolimits}
\def\Picardstack{\mathcal{P}\!\mathit{ic}}
\def\Picardfunctor{\mathrm{Pic}}
\def\Deformationcategory{\mathcal{D}\!\mathit{ef}}
\usepackage{cleveref}
\crefname{originallemma}{original lemma}{original lemmas}
\Crefname{originallemma}{Original lemma}{Original lemmas}
\makeatletter
\newlabel{lemma-DM}{{lemma-DM}{source}{}{sourcepointer-1}{}}
\newlabel{lemma-DM@cref}{{[originallemma][0][]lemma-DM}{[1][0][]source}{}{}{}}
\makeatother
\crefname{originalreference}{original reference}{original references}
\Crefname{originalreference}{Original reference}{Original references}
\makeatletter
\newlabel{spaces-descent-definition-etale-smooth-local-source-target}{{spaces-descent-definition-etale-smooth-local-source-target}{source}{}{sourcepointer-2}{}}
\newlabel{spaces-descent-definition-etale-smooth-local-source-target@cref}{{[originalreference][0][]spaces-descent-definition-etale-smooth-local-source-target}{[1][0][]source}{}{}{}}
\makeatother
\crefname{originalreference}{original reference}{original references}
\Crefname{originalreference}{Original reference}{Original references}
\makeatletter
\newlabel{spaces-descent-lemma-etale-smooth-local-source-target-characterize}{{spaces-descent-lemma-etale-smooth-local-source-target-characterize}{source}{}{sourcepointer-3}{}}
\newlabel{spaces-descent-lemma-etale-smooth-local-source-target-characterize@cref}{{[originalreference][0][]spaces-descent-lemma-etale-smooth-local-source-target-characterize}{[1][0][]source}{}{}{}}
\makeatother
\crefname{originalreference}{original reference}{original references}
\Crefname{originalreference}{Original reference}{Original references}
\makeatletter
\newlabel{spaces-descent-lemma-etale-smooth-local-source-target-implies}{{spaces-descent-lemma-etale-smooth-local-source-target-implies}{source}{}{sourcepointer-4}{}}
\newlabel{spaces-descent-lemma-etale-smooth-local-source-target-implies@cref}{{[originalreference][0][]spaces-descent-lemma-etale-smooth-local-source-target-implies}{[1][0][]source}{}{}{}}
\makeatother
\crefname{originalreference}{original reference}{original references}
\Crefname{originalreference}{Original reference}{Original references}
\makeatletter
\newlabel{stacks-properties-section-properties-morphisms}{{stacks-properties-section-properties-morphisms}{source}{}{sourcepointer-5}{}}
\newlabel{stacks-properties-section-properties-morphisms@cref}{{[originalreference][0][]stacks-properties-section-properties-morphisms}{[1][0][]source}{}{}{}}
\makeatother
\title{An etale-smooth local property on the source and target is checked on one presentation of a DM morphism}
\author{The Stacks Project Authors}
\date{stacks-project:stacks-morphisms.tex:local-snapshot-20261012; GNU FDL 1.2 or later}
\begin{document}
\maketitle
\noindent\textit{Editorial note.} All delivered mathematical text is the author's own, copied from the pinned Stacks Project chapter source sources/stacks-algebraic-stacks-morphisms/pinned-original-source.tex; no statement or proof step was rewritten, reordered or removed.\par\medskip
\noindent\textit{Editorial note.} Preamble provenance: the Stacks Project keeps one single shared preamble for every chapter; this card takes it directly from the staged sources/stacks-algebraic-stacks-morphisms/preamble.tex as three verbatim spans (lines 8-12, 14-18 and 154-217), so no derived file was written and nothing was added or redefined. The wrapper class and the base packages are supplied by the assembler.\par\medskip
\noindent\textit{Editorial note.} Bibliography provenance: neither the counted statement nor its proof contains a citation (this lane rescans the whole span before writing the card), so the card carries no bibliography and no entry was reconstructed.\par\medskip
\noindent\textit{Editorial note.} Selection provenance: this card was opened under the file-based control contract (CONTROL.md), with the mechanical pre-card gate run on its statement and proof interval before the card was written.\par\medskip
\noindent\textit{Allowed standard prerequisites (editorial).}\begin{itemize}
\item The notion of a property of morphisms of algebraic spaces being etale-smooth local on the source-and-target, in the three forms recorded in the chapter on descent on spaces. Invoked inside the counted proof (source lines 7851-7860 and 7876-7879) as cross-chapter pointers to the definition and to two lemmas. Allows the property to be moved along the smooth base change and along the etale covering used in the comparison.
\item The description of a DM morphism through the existence of an etale presentation of the fibre product over a smooth presentation of the target. Invoked at the start of the counted proof (source lines 7798-7800) as a pointer to the author lemma lemma-DM (labelled separately in the same chapter). Produces the diagram of the second condition out of the DM hypothesis.
\end{itemize}\medskip
\ifdefined\hypertarget\hypertarget{sourcepointer-1}{}\fi
\noindent\textit{Source reference.} Original lemma lemma-DM: sources/stacks-algebraic-stacks-morphisms/pinned-original-source.tex, source lines 7800--7800. Navigation pointer only: the counted text refers to the Stacks item lemma-DM, which belongs to another chapter of the project and is not part of this card. The pointer carries the target's own label as its number because this chapter's staged snapshot does not contain the other chapter's tag table; its content is not reproduced here.\par\medskip
\ifdefined\hypertarget\hypertarget{sourcepointer-2}{}\fi
\noindent\textit{Source reference.} Original reference spaces-descent-definition-etale-smooth-local-source-target: sources/stacks-algebraic-stacks-morphisms/pinned-original-source.tex, source lines 7759--7759. Navigation pointer only: the counted text refers to the Stacks item spaces-descent-definition-etale-smooth-local-source-target, which belongs to another chapter of the project and is not part of this card. The pointer carries the target's own label as its number because this chapter's staged snapshot does not contain the other chapter's tag table; its content is not reproduced here.\par\medskip
\ifdefined\hypertarget\hypertarget{sourcepointer-3}{}\fi
\noindent\textit{Source reference.} Original reference spaces-descent-lemma-etale-smooth-local-source-target-characterize: sources/stacks-algebraic-stacks-morphisms/pinned-original-source.tex, source lines 7885--7885. Navigation pointer only: the counted text refers to the Stacks item spaces-descent-lemma-etale-smooth-local-source-target-characterize, which belongs to another chapter of the project and is not part of this card. The pointer carries the target's own label as its number because this chapter's staged snapshot does not contain the other chapter's tag table; its content is not reproduced here.\par\medskip
\ifdefined\hypertarget\hypertarget{sourcepointer-4}{}\fi
\noindent\textit{Source reference.} Original reference spaces-descent-lemma-etale-smooth-local-source-target-implies: sources/stacks-algebraic-stacks-morphisms/pinned-original-source.tex, source lines 7853--7853. Navigation pointer only: the counted text refers to the Stacks item spaces-descent-lemma-etale-smooth-local-source-target-implies, which belongs to another chapter of the project and is not part of this card. The pointer carries the target's own label as its number because this chapter's staged snapshot does not contain the other chapter's tag table; its content is not reproduced here.\par\medskip
\ifdefined\hypertarget\hypertarget{sourcepointer-5}{}\fi
\noindent\textit{Source reference.} Original reference stacks-properties-section-properties-morphisms: sources/stacks-algebraic-stacks-morphisms/pinned-original-source.tex, source lines 248--248. Navigation pointer only: the counted text refers to the Stacks item stacks-properties-section-properties-morphisms, which belongs to another chapter of the project and is not part of this card. The pointer carries the target's own label as its number because this chapter's staged snapshot does not contain the other chapter's tag table; its content is not reproduced here.\par\medskip
\begin{lemma}
\label{lemma-etale-smooth-local-source-target}
Let $\mathcal{P}$ be a property of morphisms of algebraic spaces
which is \'etale-smooth local on the source-and-target.
Let $f : \mathcal{X} \to \mathcal{Y}$ be a DM morphism of algebraic stacks.
Consider commutative diagrams
$$
\xymatrix{
U \ar[d]_a \ar[r]_h & V \ar[d]^b \\
\mathcal{X} \ar[r]^f & \mathcal{Y}
}
$$
where $U$ and $V$ are algebraic spaces, $V \to \mathcal{Y}$ is smooth,
and $U \to \mathcal{X} \times_\mathcal{Y} V$ is \'etale.
The following are equivalent
\begin{enumerate}
\item for any diagram as above the morphism $h$ has property $\mathcal{P}$, and
\item for some diagram as above with $a : U \to \mathcal{X}$ surjective
the morphism $h$ has property $\mathcal{P}$.
\end{enumerate}
If $\mathcal{X}$ and $\mathcal{Y}$ are representable by algebraic spaces,
then this is also equivalent to $f$ (as a morphism of algebraic spaces)
having property $\mathcal{P}$. If $\mathcal{P}$ is also preserved under
any base change, and fppf local on the base, then for morphisms $f$
which are representable by algebraic spaces this
is also equivalent to $f$ having property $\mathcal{P}$ in the sense
of
Properties of Stacks,
Section \ref{stacks-properties-section-properties-morphisms}.
\end{lemma}

\begin{proof}
Let us prove the implication (1) $\Rightarrow$ (2). Pick an algebraic
space $V$ and a surjective and smooth morphism $V \to \mathcal{Y}$.
As $f$ is DM there exists a scheme $U$ and a surjective \'etale morphism
$U \to V \times_\mathcal{Y} \mathcal{X}$, see
Lemma \ref{lemma-DM}. Thus we see that (2) holds.
Note that $U \to \mathcal{X}$ is surjective and smooth as well, as a
composition of the base change
$\mathcal{X} \times_\mathcal{Y} V \to \mathcal{X}$ and the chosen
map $U \to \mathcal{X} \times_\mathcal{Y} V$. Hence we obtain a
diagram as in (1). Thus if (1) holds, then $h : U \to V$ has property
$\mathcal{P}$, which means that (2) holds as $U \to \mathcal{X}$ is surjective.

\medskip\noindent
Conversely, assume (2) holds and let $U, V, a, b, h$ be as in (2).
Next, let $U', V', a', b', h'$ be any diagram as in (1).
Picture
$$
\xymatrix{
U \ar[d] \ar[r]_h & V \ar[d] \\
\mathcal{X} \ar[r]^f & \mathcal{Y}
}
\quad\quad
\xymatrix{
U' \ar[d] \ar[r]_{h'} & V' \ar[d] \\
\mathcal{X} \ar[r]^f & \mathcal{Y}
}
$$
To show that (2) implies (1) we have to prove that $h'$ has $\mathcal{P}$.
To do this consider the commutative diagram
$$
\xymatrix{
U \ar[d]^h &
U \times_\mathcal{X} U' \ar[l] \ar[d]^{(h, h')} \ar[r] &
U' \ar[d]^{h'} \\
V &
V \times_\mathcal{Y} V' \ar[l] \ar[r] &
V'
}
$$
of algebraic spaces. Note that the horizontal arrows are
smooth as base changes of the smooth morphisms
$V \to \mathcal{Y}$, $V' \to \mathcal{Y}$, $U \to \mathcal{X}$, and
$U' \to \mathcal{X}$. Note that the squares
$$
\xymatrix{
U \ar[d] & U \times_\mathcal{X} U' \ar[l] \ar[d] &
U \times_\mathcal{X} U' \ar[d] \ar[r] & U' \ar[d] \\
V \times_\mathcal{Y} \mathcal{X} &
V \times_\mathcal{Y} U' \ar[l] &
U \times_\mathcal{Y} V' \ar[r] &
\mathcal{X} \times_\mathcal{Y} V'
}
$$
are cartesian, hence the vertical arrows are \'etale by our assumptions on
$U', V', a', b', h'$ and $U, V, a, b, h$.
Since $\mathcal{P}$ is smooth local on the target by
Descent on Spaces, Lemma
\ref{spaces-descent-lemma-etale-smooth-local-source-target-implies} part (2)
we see that the base change
$t : U \times_\mathcal{Y} V' \to V \times_\mathcal{Y} V'$ of $h$
has $\mathcal{P}$. Since $\mathcal{P}$ is \'etale local on the source by
Descent on Spaces, Lemma
\ref{spaces-descent-lemma-etale-smooth-local-source-target-implies} part (1)
and $s : U \times_\mathcal{X} U' \to U \times_\mathcal{Y} V'$ is \'etale,
we see the morphism $(h, h') = t \circ s$ has $\mathcal{P}$.
Consider the diagram
$$
\xymatrix{
U \times_\mathcal{X} U' \ar[r]_{(h, h')} \ar[d] &
V \times_\mathcal{Y} V' \ar[d] \\
U' \ar[r]^{h'} & V'
}
$$
The left vertical arrow is surjective, the right vertical arrow is smooth, and
the induced morphism
$$
U \times_\mathcal{X} U'
\longrightarrow
U' \times_{V'} (V \times_\mathcal{Y} V') = V \times_\mathcal{Y} U'
$$
is \'etale as seen above. Hence by
Descent on Spaces, Definition
\ref{spaces-descent-definition-etale-smooth-local-source-target} part (3)
we conclude that $h'$ has $\mathcal{P}$. This finishes the proof of the
equivalence of (1) and (2).

\medskip\noindent
If $\mathcal{X}$ and $\mathcal{Y}$ are representable, then
Descent on Spaces,
Lemma \ref{spaces-descent-lemma-etale-smooth-local-source-target-characterize}
applies which shows that (1) and (2) are equivalent to $f$ having
$\mathcal{P}$.

\medskip\noindent
Finally, suppose $f$ is representable, and $U, V, a, b, h$ are
as in part (2) of the lemma, and that $\mathcal{P}$ is preserved under
arbitrary base change. We have to show that for any scheme
$Z$ and morphism $Z \to \mathcal{X}$ the base change
$Z \times_\mathcal{Y} \mathcal{X} \to Z$
has property $\mathcal{P}$. Consider the diagram
$$
\xymatrix{
Z \times_\mathcal{Y} U \ar[d] \ar[r] &
Z \times_\mathcal{Y} V \ar[d] \\
Z \times_\mathcal{Y} \mathcal{X} \ar[r] &
Z
}
$$
Note that the top horizontal arrow is a base change of $h$ and
hence has property $\mathcal{P}$. The left vertical arrow is
surjective, the induced morphism
$$
Z \times_\mathcal{Y} U \longrightarrow
(Z \times_\mathcal{Y} \mathcal{X}) \times_Z (Z \times_\mathcal{Y} V)
$$
is \'etale, and the right vertical arrow is smooth. Thus
Descent on Spaces,
Lemma \ref{spaces-descent-lemma-etale-smooth-local-source-target-characterize}
kicks in and shows that $Z \times_\mathcal{Y} \mathcal{X} \to Z$
has property $\mathcal{P}$.
\end{proof}

\end{document}
